% GENERATED FILE. Edit canonical lessons, then run npm run generate:books.
% canonical-source-hash: fnv1a64:794853a08debf7aa
% canonical-lessons: UR-C148-chunna, UR-C148-badalna, UR-C148-udhar-lena, UR-C148-qarz-dena, UR-C148-bachana

\chapter{To choose, To change, To borrow, To lend, To save}
\label{ch:ur-to-choose-to-change-to-borrow-to-lend-to-save}
\hlchaptermodality{148}

\begin{chapteropening}
\textbf{By the end of this chapter:} \emph{I can say to choose, to change, to borrow, to lend and to save.}

Complete the last lesson of chapter 148: I can say to choose, to change, to borrow, to lend and to save.
\end{chapteropening}

\section[\texorpdfstring{\ur{چننا} (chunnā)}{چننا (chunnā)}]{\ur{چننا} (chunnā) — to choose, to pick}

\label{lesson:UR-C148-chunna}



\begin{quote}
Before the new one: say the Urdu for to smile, then the Urdu for to panic.
\end{quote}

\subsection*{You'll want to know: \ur{چننا}}
\textbf{\ur{چننا}} — \emph{chunnā} — \textquotedblleft{}to choose, to pick\textquotedblright{}.

\begin{cousinweb}
\textbf{\ur{چننا}} (\emph{chunnā}) is a tadbhava, a word inherited from Sanskrit and reshaped by sound change. It comes from \emph{cinoti}, \textquotedblleft{}heaps up, collects\textquotedblright{} (CDIAL 4814). CDIAL explains the \emph{u} as taken from other verbs.
\end{cousinweb}

\begin{grammarlens}[title={What you've built}]
One more word for this chapter.
\end{grammarlens}

\subsection*{Your turn}
\begin{itemize}
\raggedright
  \item \emph{Say it:} \emph{chunnā}
  \item \emph{Say it:} \emph{chunnā}, once more
  \item \emph{Say it:} say \emph{chunnā}
  \item \emph{Recall:} say the Urdu for to smile, then the Urdu for to panic, then say \emph{chunnā} again
\end{itemize}

\begin{tcolorbox}[breakable,colback=teal!4,colframe=teal!35!black,title={Before you move on}]
What does \ur{چننا} mean? (\textquotedblleft{}To choose, to pick\textquotedblright{}.) Say it once more.
\end{tcolorbox}
\section[\texorpdfstring{\ur{بدلنا} (badalnā)}{بدلنا (badalnā)}]{\ur{بدلنا} (badalnā) — to change}

\label{lesson:UR-C148-badalna}



\begin{quote}
Before the new one: say the Urdu for to panic, then the Urdu for to choose.
\end{quote}

\subsection*{You'll want to know: \ur{بدلنا}}
\textbf{\ur{بدلنا}} — \emph{badalnā} — \textquotedblleft{}to change\textquotedblright{}.

\begin{cousinweb}
\textbf{\ur{بدلنا}} (\emph{badalnā}) is a loan from Arabic \emph{badal}, \textquotedblleft{}a substitute, an exchange\textquotedblright{}, which reached Urdu through Persian. Its Arabic root is \emph{b-d-l} (Lane). Urdu made a verb of it with \emph{-nā}.
\end{cousinweb}

\begin{grammarlens}[title={What you've built}]
One more word for this chapter.
\end{grammarlens}

\subsection*{Your turn}
\begin{itemize}
\raggedright
  \item \emph{Say it:} \emph{badalnā}
  \item \emph{Say it:} \emph{badalnā}, once more
  \item \emph{Say it:} say \emph{badalnā}
  \item \emph{Recall:} say the Urdu for to panic, then the Urdu for to choose, then say \emph{badalnā} again
\end{itemize}

\begin{tcolorbox}[breakable,colback=teal!4,colframe=teal!35!black,title={Before you move on}]
What does \ur{بدلنا} mean? (\textquotedblleft{}To change\textquotedblright{}.) Say it once more.
\end{tcolorbox}
\section[\texorpdfstring{\ur{ادھار} \ur{لینا} (udhār lenā)}{ادھار لینا (udhār lenā)}]{\ur{ادھار} \ur{لینا} (udhār lenā) — to borrow}

\label{lesson:UR-C148-udhar-lena}



\begin{quote}
Before the new one: say the Urdu for to choose, then the Urdu for to change.
\end{quote}

\subsection*{You'll want to know: \ur{ادھار} \ur{لینا}}
\textbf{\ur{ادھار} \ur{لینا}} — \emph{udhār lenā} — \textquotedblleft{}to borrow\textquotedblright{}.

\begin{cousinweb}
\textbf{\ur{ادھار} \ur{لینا}} (\emph{udhār lenā}) joins \emph{udhār}, \textquotedblleft{}a debt, credit\textquotedblright{}, and \emph{lenā}, \textquotedblleft{}to take\textquotedblright{}. \emph{udhār} comes from Sanskrit \emph{uddhāra}, \textquotedblleft{}a raising; a loan\textquotedblright{} (CDIAL 2018), and \emph{lenā} from Sanskrit \emph{labhate}, \textquotedblleft{}catches, takes\textquotedblright{} (CDIAL 10948).
\end{cousinweb}

\begin{grammarlens}[title={What you've built}]
One more word for this chapter.
\end{grammarlens}

\subsection*{Your turn}
\begin{itemize}
\raggedright
  \item \emph{Say it:} \emph{udhār lenā}
  \item \emph{Say it:} \emph{udhār lenā}, once more
  \item \emph{Say it:} say \emph{udhār lenā}
  \item \emph{Recall:} say the Urdu for to choose, then the Urdu for to change, then say \emph{udhār lenā} again
\end{itemize}

\begin{tcolorbox}[breakable,colback=teal!4,colframe=teal!35!black,title={Before you move on}]
What does \ur{ادھار} \ur{لینا} mean? (\textquotedblleft{}To borrow\textquotedblright{}.) Say it once more.
\end{tcolorbox}
\section[\texorpdfstring{\ur{قرض} \ur{دینا} (qarz denā)}{قرض دینا (qarz denā)}]{\ur{قرض} \ur{دینا} (qarz denā) — to lend (money)}

\label{lesson:UR-C148-qarz-dena}



\begin{quote}
Before the new one: say the Urdu for to change, then the Urdu for to borrow.
\end{quote}

\subsection*{You'll want to know: \ur{قرض} \ur{دینا}}
\textbf{\ur{قرض} \ur{دینا}} — \emph{qarz denā} — \textquotedblleft{}to lend (money)\textquotedblright{}.

\begin{cousinweb}
\textbf{\ur{قرض} \ur{دینا}} (\emph{qarz denā}) joins \emph{qarz}, \textquotedblleft{}a loan\textquotedblright{}, and \emph{denā}, \textquotedblleft{}to give\textquotedblright{}. \emph{qarz} is Arabic \emph{qarḍ}, \textquotedblleft{}a loan\textquotedblright{}, which reached Urdu through Persian (Lane, root \emph{q-r-ḍ}); \emph{denā} comes from Sanskrit \emph{dadāti}, \textquotedblleft{}gives\textquotedblright{} (CDIAL 6141).
\end{cousinweb}

\begin{grammarlens}[title={What you've built}]
One more word for this chapter.
\end{grammarlens}

\subsection*{Your turn}
\begin{itemize}
\raggedright
  \item \emph{Say it:} \emph{qarz denā}
  \item \emph{Say it:} \emph{qarz denā}, once more
  \item \emph{Say it:} say \emph{qarz denā}
  \item \emph{Recall:} say the Urdu for to change, then the Urdu for to borrow, then say \emph{qarz denā} again
\end{itemize}

\begin{tcolorbox}[breakable,colback=teal!4,colframe=teal!35!black,title={Before you move on}]
What does \ur{قرض} \ur{دینا} mean? (\textquotedblleft{}To lend (money)\textquotedblright{}.) Say it once more.
\end{tcolorbox}
\section[\texorpdfstring{\ur{بچانا} (bachānā)}{بچانا (bachānā)}]{\ur{بچانا} (bachānā) — to save, to rescue}

\label{lesson:UR-C148-bachana}



\begin{quote}
Before the new one: say the Urdu for to borrow, then the Urdu for to lend.
\end{quote}

\subsection*{You'll want to know: \ur{بچانا}}
\textbf{\ur{بچانا}} — \emph{bachānā} — \textquotedblleft{}to save, to rescue\textquotedblright{}.

\begin{cousinweb}
\textbf{\ur{بچانا}} (\emph{bachānā}) is the causing form of \emph{bachnā}, \textquotedblleft{}to escape\textquotedblright{}: to save is to make someone escape. CDIAL lists \emph{bachnā} under Sanskrit \emph{vañcati}, \textquotedblleft{}moves; escapes from\textquotedblright{} (CDIAL 11208).
\end{cousinweb}

\begin{grammarlens}[title={What you've built}]
One more word for this chapter.
\end{grammarlens}

\subsection*{Your turn}
\begin{itemize}
\raggedright
  \item \emph{Say it:} \emph{bachānā}
  \item \emph{Say it:} \emph{bachānā}, once more
  \item \emph{Say it:} say \emph{bachānā}
  \item \emph{Recall:} say the Urdu for to borrow, then the Urdu for to lend, then say \emph{bachānā} again
\end{itemize}

\begin{tcolorbox}[breakable,colback=teal!4,colframe=teal!35!black,title={Before you move on}]
What does \ur{بچانا} mean? (\textquotedblleft{}To save, to rescue\textquotedblright{}.) Say it once more.
\end{tcolorbox}
